\uuid{jHNz}
\titre{Récurrence}
\auteur{bodin ; hanani ; mzari}
\organisation{exo7}
\licence{CC BY-NC-SA 4.0}
\datecreate{2020-01-20}
\niveau{L1}
\module{Algèbre}
\chapitre{Logique, ensemble, raisonnement}
\sousChapitre{Absurde et contraposée}
\difficulte{2}
\theme{Raisonnement, récurrence}
\exo7qcmid{lille-1.logique.025}

\contenu{
\begin{question}
\qclassification{100.03,100.04}
\qtags{motcle=récurrence}
Je veux montrer que $e^x > x$ pour tout $x$ réel avec $x \ge 1$.
L'initialisation est vraie pour $x=1$, car $e^1 = 2,718\ldots >1$.
Pour l'hérédité, je suppose $e^x>x$ et je calcule :
$$e^{x+1} = e^x \times e > x  \times e \ge x \times 2 \ge x + 1.$$
Je conclus par le principe de récurrence.
Pour quelles raisons cette preuve n'est pas valide ?
\begin{answers}
    \bad{Car il faudrait commencer l'initialisation à $x=0$.}

    \good{Car $x$ est un réel.}

    \bad{Car l'inégalité $e^x > x$ est fausse pour $x\le0$.}

    \bad{Car la suite d'inégalités est fausse.}
\end{answers}
\begin{explanations}
La récurrence c'est uniquement avec des entiers !
\end{explanations}
\end{question}
}
